%============ Experiencia ============%
\section*{Experiencia Profesional}

\subsection*{Contratista Independiente \hfill 2025 a la fecha}
\begin{itemize}[leftmargin=12pt]
    \item Brindo servicios de ingeniería de software e infraestructura para sistemas empresariales a medida, plataformas internas e iniciativas de habilitación técnica.
    \item Desarrollé y mantuve proyectos independientes recientes como \textit{RetailBase}, \textit{RetailBase.Identity}, \textit{poc-vite-crud}, \textit{hard.core.web}, \textit{HookHubNet}, \textit{ScepServer} y \textit{ari64}.
    \item Administré y configuré entornos VPS, conectividad VPN, servidores de correo, plataformas de colaboración y herramientas de CI/CD, incluyendo aprovisionamiento, endurecimiento, automatización de despliegues, respaldos, SSL/TLS y soporte operativo.
    \item Sigo un pipeline de desarrollo \textit{Spec-Driven} asistido de punta a punta con GitHub Copilot, que cubre descubrimiento de requerimientos, análisis, diseño de especificaciones, refinamiento, construcción de la solución, pruebas, despliegue, monitoreo y mantenimiento.
\end{itemize}

\subsection*{Desarrollador de Software / Líder Técnico \hfill Cryoinfra — 2019 a 2025}
\begin{itemize}[leftmargin=12pt]
    \item Lideré equipos multidisciplinarios basados en SCRUM, asegurando la entrega oportuna de soluciones de software de alto impacto.
    \item Diseñé, implementé y mantuve una arquitectura basada en microservicios asegurada con autenticación JWT.
    \item Incorporé herramientas de IA en el trabajo diario para acelerar corrección de código, refactorización, generación de pruebas, documentación técnica y análisis de alternativas de implementación, manteniendo validación humana y criterios de calidad de ingeniería.
    \item Comencé explorando patrones de \textit{vibe coding} en proyectos recientes para iterar con rapidez, pero evolucioné hacia un enfoque \textit{Spec-Driven} más confiable, definiendo primero requerimientos, contratos, criterios de aceptación y tareas de implementación; este enfoque se refleja especialmente en ScepServer, HookHubNet, hard.core.web y poc-vite-crud.
    \item Desarrollé aplicaciones web y móviles escalables utilizando Razor Pages, .NET Core y .NET MAUI.
    \item Integré y orquesté múltiples APIs internas y externas para la sincronización de datos entre sistemas.
    \item Definí estándares de documentación técnica, incluyendo diagramas de arquitectura, flujos de trabajo y guías de codificación.
    \item Mentoricé a desarrolladores junior, realicé revisiones de código y promoví mejores prácticas en desarrollo de software.
    \item Proporcioné soporte a usuarios y capacitación técnica para herramientas internas y aplicaciones empresariales.
\end{itemize}

\subsection*{Académico de Laboratorio \hfill Universidad Autónoma Metropolitana - Azcapotzalco — 2014 a 2017}
\begin{itemize}[leftmargin=12pt]
    \item Desarrollé y mantuve el sitio web departamental y aplicaciones usando PHP, JavaScript y MySQL.
    \item Administré servidores Linux y Windows, así como la infraestructura de red departamental.
    \item Di soporte y mantenimiento a aulas y laboratorios del Departamento de Sistemas.
    \item Contribuí a actividades de docencia e investigación apoyando a profesores y proyectos académicos.
\end{itemize}

\subsection*{Desarrollador de Software Junior \hfill Cryoinfra — 2012 a 2014}
\begin{itemize}[leftmargin=12pt]
    \item Desarrollé aplicaciones monolíticas web y de escritorio usando .NET Framework y SQL Server.
    \item Colaboré estrechamente con equipos de QA y producto para asegurar entregables de alta calidad.
    \item Diagnostiqué y resolví problemas complejos en componentes de terceros (por ejemplo, Telerik).
    \item Optimicé consultas a bases de datos, mejorando el rendimiento y la confiabilidad de las aplicaciones.
    \item Proporcioné soporte técnico y de mesa de ayuda a usuarios internos.
\end{itemize}
